% ============ font-setup.tex —— Windows 中文字体 ============
% 供 AJbook.cls 的 fontsetup 选项引入（对应 \setCJK... 与 \songti 等命令）
% AutoFakeBold/AutoFakeSlant: 中文字体无粗斜变体, 用伪粗/伪斜规避 Font Warning
\setCJKmainfont{SimSun}[AutoFakeBold=1.5, AutoFakeSlant=0.2]
\setCJKsansfont{SimHei}[AutoFakeBold=1.2, AutoFakeSlant=0.2]
\setCJKmonofont{SimSun}[AutoFakeBold=1.5, AutoFakeSlant=0.2]
\setCJKfamilyfont{song}{SimSun}[AutoFakeBold=1.5, AutoFakeSlant=0.2]
\setCJKfamilyfont{songti}{SimSun}[AutoFakeBold=1.5, AutoFakeSlant=0.2]
\setCJKfamilyfont{heiti}{SimHei}[AutoFakeBold=1.2, AutoFakeSlant=0.2]
\setCJKfamilyfont{hei2}{SimHei}[AutoFakeBold=1.2, AutoFakeSlant=0.2]
\setCJKfamilyfont{fangsong}{FangSong}[AutoFakeBold=1.5, AutoFakeSlant=0.2]
\setCJKfamilyfont{kaishu}{KaiTi}[AutoFakeBold=1.5, AutoFakeSlant=0.2]
\setCJKfamilyfont{sectionfont}{SimHei}[AutoFakeBold=1.2, AutoFakeSlant=0.2]
\setCJKfamilyfont{pffont}{KaiTi}[AutoFakeBold=1.5, AutoFakeSlant=0.2]
% 西里尔字符（如 Чеботарёв）用 CJK 字体渲染（SimSun 含西里尔字形）
\xeCJKDeclareCharClass{CJK}{"0400 -> "04FF}
% 保留 CJK 字符间的显式空格（西里尔词与中文之间留空隙）
\xeCJKsetup{CJKspace=true}
\newcommand{\songti}{\CJKfamily{songti}}
\newcommand{\heiti}{\CJKfamily{heiti}}
\newcommand{\fangsong}{\CJKfamily{fangsong}}
\newcommand{\kaishu}{\CJKfamily{kaishu}}
\newcommand{\thmheiti}{\CJKfamily{heiti}}
